problem:  
Let \((X,d,\mathfrak{m})\) be a complete, essentially non-branching metric measure space satisfying the curvature–dimension condition \(CD(K,N)\) for some \(K\in\mathbb{R}\), \(1<N<\infty\).  
Assume that for each \(x\in X\), the measured tangent cone \(T_xX\) (in the pointed measured Gromov–Hausdorff sense) is uniquely Euclidean, i.e.  
\[
T_xX \cong (\mathbb{R}^{k(x)},|\cdot|,\lambda^{k(x)}),
\]
and that the dimension function \(x\mapsto k(x)\) is upper semicontinuous and locally constant on a dense open subset of \(X\).  

**Question:**  
Must \(k(x)\) be constant across all of \(X\), and must the Cheeger energy \(\mathrm{Ch}\) on \(L^2(X,\mathfrak{m})\) be quadratic—so that \((X,d,\mathfrak{m})\) is infinitesimally Hilbertian, hence an \(RCD(K,N)\) space?  
Equivalently: can a \(CD(K,N)\) space have Euclidean tangent cones of locally constant dimension without being Hilbertian—i.e. can synthetic Ricci curvature allow “locally Riemannian but globally non‑Riemannian” phenomena?  

Why is it a "good" problem:  
This problem targets the **rigidity of Euclidean tangent structure** within synthetic lower Ricci curvature geometry. Known examples of non‑Hilbertian \(CD(K,N)\) spaces arise from **Finsler geometries**, where all tangents are **Minkowski spaces**—strictly convex normed vector spaces that are non‑Euclidean. Hence, by explicitly assuming *all* tangents are Euclidean, the standard Finsler obstruction to Hilbertianity is removed. The question then isolates a deeper possibility: even with genuinely Euclidean infinitesimal data, can non‑Hilbertian analytic behavior persist?  

A positive answer (that Euclidean tangents and local dimension regularity force \(RCD\)) would establish a strong **dimension‑propagation and Hilbertian‑rigidity principle** under \(CD(K,N)\). A negative answer would reveal entirely new types of \(CD\) spaces—neither Riemannian nor classically Finsler—forcing a reevaluation of how synthetic curvature interacts with measure and first‑order structure.  

The hypothesis of upper semicontinuity and dense open constancy is the minimal setting that allows the possibility of dimension jumps while retaining local geometric coherence. Resolving whether these weak regularity conditions suffice to enforce global constancy links subtle themes in **measured Gromov–Hausdorff geometry**, **geometric measure theory**, and **Dirichlet form analysis**.  

The problem is conceptually simple (“do Euclidean tangents with locally constant dimension force global Hilbertianity?”), yet requires fundamental advances across multiple subfields to answer—embodying the “narrow entrance, profound depth” criterion of a good problem.